\section{Design decisions and performance}
\label{sec:optimization}

This section records the design decisions that shape the shipped
encoder's cost --- both wall-clock and proof size --- and the
reasoning behind each. The single most consequential one is the
radix-8 evaluator; the rest are coefficient-representation and
parallelism choices.

\subsection{Radix-8 evaluator: the coefficient-size win}
\label{sec:opt-radix8}

The dominant design choice is the radix-8 evaluator
(\autoref{sec:radix8}). Its purpose is \emph{not} cache locality ---
it is to control the magnitude of the lifted codeword coefficients.

A codeword cell is $P(v_k) = \sum_i c_i\, X_i(v_k)$ evaluated in
$\Rlift$. How that sum is computed determines how large the
$\Z$-coefficients of the result grow:

\begin{itemize}
  \item \textbf{Radix-2 $\Z$-lift.} The FFT is a chain of $m$
    sequential multiplications in $\Rlift$ by lifted twiddles. Each is
    a $\Z$-arithmetic \code{mul + reduce-mod-}$\flift$; $\Z$-carries
    accumulate across all $m$ stages. Empirically the codeword
    coefficients reach $\sim 2^{47}$ ($49$ signed bits) at
    $\code{codeword\_len} = 4096$ --- they do not fit in \code{i32}
    and need \code{i64} storage and a $\ge 49$-bit wire encoding.
  \item \textbf{Radix-8 with $\GF$-precomputed Vandermonde.} Each
    $8 \times 8$ block matrix is computed once in $\GF(2^{16})$ --- so
    every entry is a $\{0,1\}$-coefficient lifted element --- and the
    FFT composes only $m/3$ such matvecs in $\Rlift$. With a third as
    many carry-accumulating meta-stages, and a $\{0,1\}$-coefficient
    multiplier at each, the codeword coefficients stay near $2^{26}$
    ($28$ signed bits) at the same $\code{codeword\_len} = 4096$.
\end{itemize}

That is a $\sim 21$-bit reduction in coefficient magnitude. It is
\emph{not} free: the radix-8 form does more arithmetic overall
(an $8\times 8$ matvec is $64$ multiplications per block of $8$, versus
$\sim 8$ for three radix-2 stages), so commit time is somewhat higher
than a radix-2 lift would be. The trade is deliberate: the bit-size win
is what lets each codeword cell pack into $32$ bits on the wire
(\autoref{sec:opt-packing}), which is the dominant term of the proof
size. \autoref{sec:measurements} quantifies both sides.

\paragraph{It is a different code.} As noted in \autoref{sec:radix8},
the radix-8-with-$\GF$-lifted-Vandermonde encoder produces a different
lifted code than a radix-2 $\Z$-lift --- the two agree mod $2$ but not
in the higher $\Z$-bits. The small-coefficient variant is the one we
ship.

\subsection{Native {\tt i64} coefficient storage}
\label{sec:opt-native}

Codeword and combined-row cells store their $16$ coefficients as
native \code{i64} (\code{DensePolynomial<i64, 16>}), not
\code{crypto-bigint}'s multi-limb \code{Int<N>}. The radix-8
coefficient bound ($\sim 2^{26}$ at our workload, growing toward
$2^{63}$ only at the largest $\code{codeword\_len} = 2^{16}$) fits
\code{i64}, so the FFT runs on single-instruction $64$-bit ALU adds
instead of generic limb-by-limb carry propagation. The
\code{FftRingElement16} trait keeps the kernel generic, so an
\code{Int<N>} instantiation remains available if a future workload
needs more than $64$ bits per coefficient.

\subsection{Bit-packed wire encoding}
\label{sec:opt-packing}

In-memory width and on-wire width are decoupled. Cells are
\code{i64}-backed in memory (\autoref{sec:opt-native}) but the
\code{Transcribable} encoding of \code{PackedI64Poly16<BITS>}
(\autoref{sec:packed-cw}) writes only \code{BITS} bits per
coefficient.

\begin{itemize}
  \item \textbf{Codeword cells: \code{BITS\_CW = 32}.} The radix-8
    codeword coefficients max near $2^{26}$, so $32$ bits is lossless
    with $\sim 6$ bits of headroom. Column openings are the dominant
    proof component, so this width sets most of the proof size.
  \item \textbf{Combined-row cells: \code{BITS\_CR = 8}.} At
    $\code{num\_rows} = 1$ a combined-row coefficient is bounded by
    $\code{batch}\cdot 15$ (\autoref{sec:protocol-integration}), which
    fits $8$ signed bits at $\code{batch} = 11$. Packing the
    combined row at $8$ bits rather than the full $64$ shrinks that
    proof component $\sim 8\times$.
\end{itemize}

Splitting the two widths matters because the codeword and combined-row
cells have very different magnitude budgets --- codeword coefficients
are FFT-amplified, combined-row coefficients are not.

\subsection{Batched $\flift$-reduction in the matvec}
\label{sec:opt-batched-reduce}

Each $8 \times 8$ block matvec produces one output cell as a sum of
eight lifted-twiddle products. Reducing mod $\flift$ after every one
of the eight multiplications wastes work: the
\code{FftRingElement16::UnreducedAcc} path
(\autoref{sec:fft-ring-element}) accumulates all eight products in a
single length-$31$ unreduced buffer and folds it back mod $\flift$
exactly once per output cell --- one reduction instead of eight.

\subsection{The {\tt parallel} feature}
\label{sec:opt-parallel}

Under the \code{parallel} feature the per-meta-stage block loop is a
rayon \code{par\_chunks\_mut} (via \code{cfg\_chunks\_mut!}). Each
block of $8$ is a disjoint slice, so block-level parallelism is
data-race-free. At commit time there are also \code{batch} independent
row FFTs, giving rayon ample work to distribute. This is the largest
single wall-clock lever; benchmarks in \autoref{sec:measurements} are
all run with \code{parallel} enabled.

\subsection{Zero-twiddle skip}

Vandermonde entries equal to $0 \in \GF(2^{16})$ contribute nothing
and are skipped in the matvec inner loop; the $X_0 \equiv 1$ column,
whose entries are all $1$, contributes a plain add with no
lifted-twiddle multiplication. Small but free.
